\chapter{Limits and Continuity}\label{ch:tle:limcont}

This chapter extends the notion of limit from sequences to functions of a
real variable, and introduces continuity. The central result is the
intermediate value theorem, which turns qualitative information (``$f$ is
continuous and changes sign'') into the existence of solutions of equations.

\section{Limits of functions}

Throughout, $f$ is a function defined on an interval or a union of intervals
$D \subseteq \R$.

\begin{definition}[Limit at infinity]\label{def:tle:limcont:liminf}
Let $f$ be defined on an interval $\intco{a}{+\infty}$ and $\ell \in \R$. We
say that $f$ \emph{tends to $\ell$ at $+\infty$}\index{limit!of a function},
written $\lim\limits_{x\to+\infty} f(x) = \ell$, if every open interval
containing $\ell$ contains all values $f(x)$ for $x$ large enough:
for every $\varepsilon > 0$ there exists $A$ such that for all $x \geq A$,
$\abs{f(x) - \ell} \leq \varepsilon$.

We say $f$ tends to $+\infty$ at $+\infty$ if for every $M \in \R$ there
exists $A$ such that $f(x) \geq M$ for all $x \geq A$. Limits at $-\infty$
and limits equal to $-\infty$ are defined analogously.
\end{definition}

\begin{definition}[Limit at a point]\label{def:tle:limcont:limpoint}
Let $a \in \R$ and let $f$ be defined near $a$ (except possibly at $a$). We
say that $f$ tends to $\ell \in \R$ at $a$, written
$\lim\limits_{x \to a} f(x) = \ell$, if for every $\varepsilon > 0$ there
exists $\delta > 0$ such that for all $x \in D$ with
$\abs{x - a} \leq \delta$, $\abs{f(x) - \ell} \leq \varepsilon$.

We say $f$ tends to $+\infty$ at $a$ if for every $M$ there is $\delta > 0$
such that $f(x) \geq M$ whenever $x \in D$ and $\abs{x - a} \leq \delta$.
One-sided limits ($x \to a^+$, $x \to a^-$) restrict $x$ to $x > a$ or
$x < a$.
\end{definition}

\begin{example}
$\lim\limits_{x\to+\infty} \frac{1}{x} = 0$,
$\lim\limits_{x\to 0^+} \frac{1}{x} = +\infty$,
$\lim\limits_{x\to 0^-} \frac{1}{x} = -\infty$. The function
$x \mapsto \frac1x$ has no (two-sided) limit at $0$.
\end{example}

\begin{definition}[Asymptotes]\label{def:tle:limcont:asymptote}
The line $y = \ell$ is a \emph{horizontal asymptote}\index{asymptote} of the
curve of $f$ at $+\infty$ (resp.\ $-\infty$) if
$\lim_{x\to+\infty} f(x) = \ell$ (resp.\ at $-\infty$). The line $x = a$ is a
\emph{vertical asymptote} if $f$ has an infinite one-sided limit at $a$.
\end{definition}

\begin{proposition}[Operations on limits]\label{prop:tle:limcont:operations}
The rules of \cref{prop:tle:seq:operations} (sums, products, quotients) hold
verbatim for limits of functions, at a point or at infinity, with the same
four indeterminate forms $(+\infty)+(-\infty)$, $0\times\infty$,
$\frac{\infty}{\infty}$, $\frac{0}{0}$.
\end{proposition}

\begin{proof}
The proofs are identical to the sequence case, replacing ``for $n \geq N$''
by ``for $x \geq A$'' or ``for $\abs{x-a} \leq \delta$''.
\end{proof}

\begin{theorem}[Limit of a composition]\label{thm:tle:limcont:composition}
Let $f, g$ be functions and $a, b, c$ be real numbers or $\pm\infty$. If
$\lim\limits_{x \to a} f(x) = b$ and $\lim\limits_{y \to b} g(y) = c$, then
\[
\lim_{x \to a} g\bigl(f(x)\bigr) = c .
\]
The same statement holds with a sequence in place of $f$: if
$u_n \to b$ and $\lim_{y\to b} g(y) = c$, then $g(u_n) \to c$.
\end{theorem}

\begin{proof}
Take $b, c$ finite (the other cases are analogous). Let $\varepsilon > 0$.
There is $\delta > 0$ such that $\abs{y - b} \leq \delta$ implies
$\abs{g(y) - c} \leq \varepsilon$; and there is $\delta' > 0$ such that
$\abs{x - a} \leq \delta'$ implies $\abs{f(x) - b} \leq \delta$. Chaining the
two gives $\abs{g(f(x)) - c} \leq \varepsilon$ for
$\abs{x - a} \leq \delta'$.
\end{proof}

\begin{theorem}[Comparison and squeeze]\label{thm:tle:limcont:squeeze}
Let $f, g, h$ be functions and $a \in \R \cup \{\pm\infty\}$.
\begin{enumerate}
  \item If $f \leq g$ near $a$ and $f(x) \to +\infty$ as $x \to a$, then
        $g(x) \to +\infty$.
  \item If $f \leq g \leq h$ near $a$ and $f, h$ tend to the same finite
        limit $\ell$ at $a$, then $g(x) \to \ell$ as $x \to a$.
\end{enumerate}
\end{theorem}

\begin{proof}
Same argument as for sequences (\cref{thm:tle:seq:squeeze}).
\end{proof}

\begin{method}[Computing limits in practice]\label{met:tle:limcont:limits}
\begin{itemize}
  \item Polynomials and rational functions at $\pm\infty$: factor out the
        highest power of $x$; the limit is that of the ratio of leading
        terms.
  \item $\frac{0}{0}$ at a point $a$: factor $(x-a)$ out of numerator and
        denominator, or recognize a derivative
        $\lim_{x\to a}\frac{f(x)-f(a)}{x-a} = f'(a)$ (\cref{ch:tle:deriv}).
  \item Square roots: multiply by the conjugate expression.
  \item Oscillating factors ($\cos$, $\sin$): use the squeeze theorem.
\end{itemize}
\end{method}

\section{Continuity}

\begin{definition}[Continuity]\label{def:tle:limcont:continuity}
Let $f$ be defined on an interval $I$ and $a \in I$. The function $f$ is
\emph{continuous at $a$}\index{continuity} if
$\lim\limits_{x \to a} f(x) = f(a)$. It is \emph{continuous on $I$} if it is
continuous at every point of $I$.
\end{definition}

Intuitively, the graph of a continuous function on an interval can be drawn
without lifting the pen.

\begin{proposition}[Continuity of usual functions]\label{prop:tle:limcont:usual}
Polynomials, rational functions, $x \mapsto \sqrt{x}$, $x \mapsto \abs{x}$,
$\cos$, $\sin$, $\exp$ and $\ln$ are continuous on every interval contained
in their domain. Sums, products, quotients (where defined) and compositions
of continuous functions are continuous.
\end{proposition}

\begin{proof}[Partial proof]
Stability under sums, products, quotients and composition follows from the
corresponding theorems on limits. Continuity of $x \mapsto x$ and of constant
functions is immediate from the definition, hence polynomials and rational
functions are continuous. Continuity of the remaining usual functions is
proved in the chapters where they are constructed, or admitted.
\end{proof}

\begin{example}
The \emph{floor function} $x \mapsto \floor{x}$ (the greatest integer
$\leq x$) is not continuous at any integer $n$: its left limit there is
$n - 1$ while $\floor{n} = n$.
\end{example}

\begin{proposition}[Sequences and continuity]\label{prop:tle:limcont:seqcont}
If $u_n \to \ell$ and $f$ is continuous at $\ell$, then $f(u_n) \to f(\ell)$.
In particular, if a sequence defined by $u_{n+1} = f(u_n)$ converges to
$\ell$ and $f$ is continuous at $\ell$, then $f(\ell) = \ell$.
\end{proposition}

\begin{proof}
The first claim is \cref{thm:tle:limcont:composition} applied with the
sequence $(u_n)$. For the second, pass to the limit on both sides of
$u_{n+1} = f(u_n)$: the left side tends to $\ell$, the right side to
$f(\ell)$, and limits are unique.
\end{proof}

\section{The intermediate value theorem}

\begin{theorem}[Intermediate value theorem]\label{thm:tle:limcont:ivt}
\index{intermediate value theorem}
Let $f$ be continuous on $\intcc{a}{b}$ and let $k$ be any number between
$f(a)$ and $f(b)$. Then there exists $c \in \intcc{a}{b}$ such that
$f(c) = k$.
\end{theorem}

\begin{proof}
Replacing $f$ by $f - k$ (and by $k - f$ if needed), we may assume
$f(a) \leq 0 \leq f(b)$ and look for a zero of $f$. We construct two
sequences by \emph{dichotomy}\index{dichotomy}: set $a_0 = a$, $b_0 = b$;
given $a_n \leq b_n$ with $f(a_n) \leq 0 \leq f(b_n)$, let
$m = \frac{a_n + b_n}{2}$ and define
\[
(a_{n+1}, b_{n+1}) =
\begin{cases}
(a_n, m) & \text{if } f(m) \geq 0,\\
(m, b_n) & \text{if } f(m) < 0.
\end{cases}
\]
In both cases $f(a_{n+1}) \leq 0 \leq f(b_{n+1})$, the sequence $(a_n)$ is
increasing, $(b_n)$ is decreasing, and
$b_n - a_n = \frac{b-a}{2^n} \to 0$: the sequences are adjacent (see
\cref{exo:tle:seq:10}) and converge to a common limit
$c \in \intcc{a}{b}$. By continuity, $f(a_n) \to f(c)$ and
$f(b_n) \to f(c)$; passing to the limit in $f(a_n) \leq 0$ and
$f(b_n) \geq 0$ gives $f(c) \leq 0 \leq f(c)$, so $f(c) = 0$.
\end{proof}

\begin{remark}
The dichotomy proof is an algorithm: it yields arbitrarily precise
approximations of a solution, halving the error at each step. This is how a
computer can solve $f(x) = 0$ knowing only how to evaluate $f$.
\end{remark}

\begin{theorem}[Bijection theorem]\label{thm:tle:limcont:bijection}
Let $f$ be continuous and \emph{strictly monotonic} on an interval $I$. Then
for every $k$ strictly between the limits of $f$ at the endpoints of $I$, the
equation $f(x) = k$ has a \emph{unique} solution in $I$.
\end{theorem}

\begin{proof}
Existence: choose $a, b \in I$ with $k$ between $f(a)$ and $f(b)$ (possible
since the values of $f$ approach its endpoint limits), and apply the
intermediate value theorem on $\intcc{a}{b}$. Uniqueness: if $x < y$, strict
monotonicity gives $f(x) \neq f(y)$, so two distinct points cannot both be
solutions.
\end{proof}

\begin{method}[Solving $f(x) = k$ with the bijection theorem]\label{met:tle:limcont:bijection}
To prove that $f(x) = k$ has exactly one solution in an interval:
\begin{enumerate}
  \item establish a variation table of $f$ (sign of $f'$, see
        \cref{ch:tle:deriv});
  \item on each interval where $f$ is continuous and strictly monotonic,
        compare $k$ with the endpoint values or limits;
  \item apply the bijection theorem on each such interval and count the
        solutions.
\end{enumerate}
A calculator or dichotomy then localizes each solution as finely as desired.
\end{method}

\begin{example}\label{ex:tle:limcont:cubic}
The equation $x^3 + x = 1$ has a unique real solution. Indeed
$f(x) = x^3 + x$ is continuous and strictly increasing on $\R$ (as a sum of
strictly increasing functions), $\lim_{-\infty} f = -\infty$ and
$\lim_{+\infty} f = +\infty$, so the bijection theorem applies with $k = 1$.
Since $f(0.6) = 0.816 < 1$ and $f(0.7) = 1.043 > 1$, the solution lies in
$\intoo{0.6}{0.7}$.
\end{example}

\section{Exercises}

\begin{exercise}[$\star$]\label{exo:tle:limcont:1}
Compute the following limits:
\[
\lim_{x\to+\infty} \frac{2x^3 - x + 1}{5x^3 + x^2}, \qquad
\lim_{x\to-\infty} \bigl(x^2 + 3x - 1\bigr), \qquad
\lim_{x\to 2^+} \frac{x+1}{x-2}, \qquad
\lim_{x\to+\infty} \frac{\cos x}{x}.
\]
\end{exercise}

\begin{exercise}[$\star$]\label{exo:tle:limcont:2}
Let $f(x) = \dfrac{2x^2 - 3x + 1}{x - 1}$, defined for $x \neq 1$.
\begin{enumerate}
  \item Show that $f(x) = 2x - 1$ for all $x \neq 1$.
  \item Deduce $\lim\limits_{x \to 1} f(x)$. Can $f$ be extended to a
        function continuous at $1$?
\end{enumerate}
\end{exercise}

\begin{exercise}[$\star$]\label{exo:tle:limcont:3}
Determine the asymptotes of the curve of
$f(x) = \dfrac{3x - 1}{x + 2}$ and sketch the curve.
\end{exercise}

\begin{exercise}[$\star\star$]\label{exo:tle:limcont:4}
Compute
\[
\lim_{x\to+\infty} \bigl(\sqrt{x^2 + x} - x\bigr)
\qquad\text{and}\qquad
\lim_{x\to 0} \frac{\sqrt{1+x} - 1}{x}.
\]
\end{exercise}

\begin{exercise}[$\star\star$]\label{exo:tle:limcont:5}
Let $f(x) = x + \sin x$.
\begin{enumerate}
  \item Show that $f(x) \to +\infty$ as $x \to +\infty$, although $f$ is not
        monotonic on any interval $\intco{A}{+\infty}$.
  \item Show that $g(x) = \dfrac{x + \sin x}{x - \sin x}$ tends to $1$ at
        $+\infty$.
\end{enumerate}
\end{exercise}

\begin{exercise}[$\star\star$]\label{exo:tle:limcont:6}
Show that the equation $x^3 - 3x + 1 = 0$ has exactly three real solutions,
and give for each an interval of length $1$ containing it.
(Hint: study the variations of $f(x) = x^3 - 3x + 1$; its derivative is
$3x^2 - 3$.)
\end{exercise}

\begin{exercise}[$\star\star$]\label{exo:tle:limcont:7}
Let $f \colon \intcc{0}{1} \to \intcc{0}{1}$ be continuous. Show that $f$ has
a \emph{fixed point}: there exists $c \in \intcc{0}{1}$ with $f(c) = c$.
(Hint: apply the intermediate value theorem to $g(x) = f(x) - x$.)
\end{exercise}

\begin{exercise}[$\star\star\star$]\label{exo:tle:limcont:8}
A hiker climbs a mountain trail, starting at 8:00 and arriving at 20:00. The
next day she descends the same trail, again starting at 8:00 and arriving at
20:00. Show that there is a time of day at which she is at exactly the same
point of the trail on both days. (Hint: introduce the difference of the two
position functions.)
\end{exercise}

\begin{exercise}[$\star\star\star$]\label{exo:tle:limcont:9}
Let $f(x) = x^5 + x^3 + x - 1$.
\begin{enumerate}
  \item Show that $f$ has a unique real zero $\alpha$, and that
        $\alpha \in \intoo{0}{1}$.
  \item Describe a dichotomy procedure computing $\alpha$ to within
        $10^{-6}$, and give the number of iterations required starting from
        $\intcc{0}{1}$.
\end{enumerate}
\end{exercise}
